\section{Modality coverage map}
\label{sec:coverage}
Table~\ref{tab:coverage} and Figure~\ref{fig:coverage} summarise, for each modality family,
the best public pairing of innate immune cell imaging with that modality and what it
supports. Three things stand out. First, same-cell pairings exist only for
immunophenotype (morphology classes, cell type), metabolism (NAD(P)H lifetime) and
signalling (reporter dynamics), and none of them involves a memory state; the one
memory-adjacent same-cell resource, sequential NF-\(\kappa\)B stimulation, has two wells per
arm. Second, transcriptome and chromatin data are paired with dynamics only at the level of
five conditions, which supports description, not prediction. Third, the families in which
innate memory is actually measured, cytokines after restimulation and the transcriptional
and epigenetic programmes behind them, have no public imaging pairing at all; the
public restimulation transcriptomes (Section~\ref{sec:partC}) are what an imaging study of
memory would have to be validated against. Proteomics beyond signalling reporters,
integrated single-cell multi-omics, spatial omics, perturbation screens, repertoire,
microbiome and clinical data have no public pairing with in-vitro innate cell imaging.
For the same reason we did not project image features onto molecular activation axes: no
public dataset has both images and a stimulus-matched molecular anchor.

Within the same-cell pairings, frozen image or video features carry information about
morphology classes and NF-\(\kappa\)B dynamics carry information about stimulus identity, but
where a simple measured quantity is available as a reference (autofluorescence intensity
for cell identity; the photon-count image for lifetime), the learned features did not add
to it.
